% Generated draft: requires mathematical and source review.
\documentclass[11pt]{article}
\usepackage[margin=25mm]{geometry}
\usepackage{fontspec}
\setmainfont{Latin Modern Roman}
\usepackage{amsmath,amssymb,mathtools,unicode-math}
\setmathfont{Latin Modern Math}
\usepackage{enumitem,needspace,xurl,hyperref}
\setlength{\parindent}{0pt}
\setlength{\parskip}{5pt}
\setlength{\emergencystretch}{3em}
\newcommand{\N}{\mathbb{N}}
\newcommand{\Z}{\mathbb{Z}}
\newcommand{\Q}{\mathbb{Q}}
\newcommand{\R}{\mathbb{R}}
\newcommand{\C}{\mathbb{C}}
\newcommand{\F}{\mathbb{F}}
\newcommand{\A}{\mathbb{A}}
\newcommand{\PP}{\mathbb{P}}
\newcommand{\E}{\mathbb{E}}
\newcommand{\eps}{\varepsilon}
\newcommand{\dd}{\,\mathrm{d}}
\begin{document}
\section*{UnsolvedMath ID: 20002679}
\textbf{Source problem code:} catalogue entry
\subsection{Definitions and mathematical statement}
The supplied verification assigns Question 19 to the unbounded Haagerup--Schultz extension problem. The extraction also contains an independent Ozawa conjecture and the opening of a random-matrix discussion. All are retained separately; the string 40.7 is reported to combine page number 4 and section 0.7.

\textbf{Primary question.} In a finite tracial von Neumann algebra \((M,\tau)\), write \(\widetilde M\) for the algebra of affiliated measurable operators and define
\[
 L^p(M,\tau)=\{T:\tau(|T|^p)<\infty\}\quad(p>0),\qquad
 M^\Delta=L_{\log}(M,\tau)
       =\{T:\tau(\log(1+|T|))<\infty\}.
\]
Does the main Haagerup--Schultz result extend from bounded \(T\in M\) to \(T\in L^p(M)\), for some or all \(p\), to all \(M^\Delta\), or to all \(\widetilde M\)? The source does not specify which part of the original result is meant. A strong interpretation asks for original-algebra, unique, hyperinvariant Brown-splitting projections: for Borel \(E\subset\C\), a projection \(p\in M\) invariant under \(T\), with \(\tau(p)=\mu_T(E)\), and Brown measures of the two diagonal compressions supported on \(E\) and its complement, respectively, with the appropriate normalized restrictions when their traces are nonzero. Hyperinvariance refers to invariance under the bounded commutant of \(T\), with the necessary domain conventions for unbounded operators. Brown measure \(\mu_T\) in its usual logarithmic-integrability framework is available for \(T\in M^\Delta\); no unrestricted Brown measure is specified for all affiliated operators.

The aid reports
\[
 L^p(M)\subset M^\Delta\quad(p>0),
\]
and attributes to Dykema, Sukochev and Zanin a Brown-splitting theorem in an enlarged finite algebra containing \(M\), applicable to all \(M^\Delta\) and hence every positive \(p\). An enlarged-algebra splitting projection is not thereby a projection in \(M\), and does not establish the full original-algebra uniqueness/hyperinvariance package. The latter general extension remains unverified in the supplied attempt.

The reported class boundaries are strict in every diffuse finite algebra:
\[
 \bigcup_{p>0}L^p(M)\ \subsetneq\ M^\Delta\ \subsetneq\ \widetilde M.
\]
In a diffuse abelian copy of \(L^\infty(0,1)\), illustrative affiliated multiplication functions are \(T_1(s)=e^{1/\sqrt s}\), which is log integrable but belongs to no positive \(L^p\), and \(T_2(s)=e^{1/s}\), which is measurable but not log integrable. For normal affiliated \(T\), the aid reports original-algebra hyperinvariant spectral projections; when \(T\in M^\Delta\), these give Brown-splitting projections. The general nonnormal case cannot be transferred by assuming modulus cutoffs are invariant. For example,
\[
 T=\begin{pmatrix}0&2\\1&0\end{pmatrix},\quad
 |T|=\begin{pmatrix}1&0\\0&2\end{pmatrix},\quad
 p=\mathbf1_{[0,3/2]}(|T|)=\begin{pmatrix}1&0\\0&0\end{pmatrix}
\]
has \((1-p)Tp\ne0\), so even this bounded cutoff subspace is not \(T\)-invariant.

\textbf{Appended Ozawa conjecture, a separate question.} Let \(M=L(\mathbb F_r)\) be a free group factor, and let \(H\) be a nonzero \(M\)-\(M\) Hilbert correspondence weakly contained in the coarse correspondence \(C=L^2(M)\otimes L^2(M)\). The appended conjecture asks whether
\[
 \operatorname{Hom}_{M\text{-}M}
       (H\otimes_MH\otimes_MH,\ C)\ne\{0\}.
\]
Here the tensor products are Connes relative tensor products and \(\operatorname{Hom}\) means bounded bimodule intertwiners. Nonzero \(H\) is necessary and is not explicit in the extraction. The weak-containment condition is expressed there by continuity of
\[
 M\odot M^{\mathrm{op}}\longrightarrow B(H),\qquad
 x\otimes y^{\mathrm{op}}\longmapsto(\xi\mapsto x\xi y)
\]
for the minimal tensor norm; the opposite algebra makes the right action precise. The source gives examples from the basic construction over a hyperfinite \(A\subset M\), interpreted as the completed correspondence \(L^2(M)\otimes_A L^2(M)\). This conjecture is not addressed by the aid's affiliated-operator results. The source does not specify the free-group rank range or all correspondence conventions.

\textbf{Appended Itzykson--Zuber context.} The next section introduces random matrix laws \(\mu_{A_N},\mu_{B_N}\) and prints
\[
 \operatorname{IZ}(A_N,B_N)=\iint
 e^{-N\operatorname{Tr}(AU^*BU)}\,
                  \dd\mu_{A_N}(A)\,\dd\mu_{B_N}(B).
\]
It attributes to Guionnet and Zeitouni an asymptotic \(\operatorname{IZ}\sim e^{-N\psi}\) under uniform operator-norm bounds. The displayed integral leaves \(U\) and its averaging unspecified, and omits spectral convergence hypotheses and trace/exponential normalizations. It is retained as incomplete reported context, not silently changed into a standard Haar integral or an asserted asymptotic theorem. No further question follows within the supplied record.
\subsection{Short English statement}
The assigned question asks how far Haagerup--Schultz projections extend to unbounded operators. The aid separates log-integrable enlarged-algebra splitting from the stronger original-algebra package and identifies sharp integrability boundaries. The extraction also contains an independent free-group-factor conjecture and incomplete Itzykson--Zuber context, which are preserved without claiming they are resolved.
\subsection{Sources}
\begin{itemize}
\item Local UnsolvedMath registry, permanent ID 20002679, catalogue entry; canonical \texttt{aim-probability-notes.json}, notes[120], reconstructed-unverified source status.
\item AIM \emph{Free Analysis}, end of section 0.6, section 0.7 ``Free Group Factors'' (Ozawa), and opening of section 0.8 ``Focus Group on Combinatorics of Random Matrix Models (day 4)'', \url{https://aimath.org/WWN/freeanalysis/freeanalysis.pdf}. The supplied verification identifies Question 19 as the assigned primary question and explains the extraction fusion.
\item Source attributions: Haagerup and Schultz; Ozawa; Guionnet and Zeitouni. Aid attribution for the enlarged-algebra theorem: Dykema, Sukochev and Zanin. No full bibliographic entries or theorem numbers are supplied.
\item Machine-generated research attempt 1, partial-result/partially-solved labels, reported rigor proved and low novelty confidence. Its extension ladder and obstructions are not peer reviewed and require independent verification before citation.
\end{itemize}
\end{document}
